The analysis of experimental data was handled by CabanaBoy. 
Pool depth of 30 was used for combinatorial background invariant mass ($M_{\text{inv}}$) distribution (BG). 
The data was divided into 30 bins (3 vertex and 10 centrality). 
The tail (10\% of statistics from the right edge of the distribution) of foreground $M_{\text{inv}}$ distribution (FG) has very low non-combinatorial contamination therefore for each centrality, vertex, and $p_T$ bin the BG was scaled so that it's tail integral matches the tail integral of FG.
Additional re-scaling factor of 0.95 for BG was introduced to avoid overestimation of BG scaling due to presence of low non-combinatorial contamination and in some cases insufficient statistics.
BG was then subtracted from the FG and obtained distributions were merged for all vertex bins and for some centrality bins according to centrality measurement ranges.

Obtained $M_{\text{inv}}$ distributions for different centrailty and $p_T$ ranges were approximated with the sum of background function and convolution of gauss and Relativistic Breit-Wigner (RBW) function.
For background function 3rd order polynomial was used for \nopid and 2nd order polynomial was used for \twopid and \onepid.
The convolution with gauss is needed to account for gaussian broadening $\sigma_{\text{gb}}$ of \kstar signal (Sec. \ref{sec:analysis_gaussian_broadening}).
RBW is presented below.

\begin{equation}
\text{RBW}(M_{\text{inv}}) = \frac{A}{2\pi} \frac{M_{\text{inv}} M_{\text{mean}} \Gamma}{\left( M_{\text{inv}}^2 - M_{\text{mean}}^2 \right)^2 + M_{\text{mean}}^2 \Gamma^2}
\end{equation}

Where $A$ - scale parameter, $M_{\text{mean}}$ and $\Gamma$ - mean and width of \kstar RBW respectively. 
For each approximation $M_{\text{mean}}$ was limited within 5\% of PDG value (891.66 GeV/$c^2$), $\Gamma$ was limited within 10\% of PDG value (50.8 GeV/$c^2$), width of gauss was limited within 10 \% of $\sigma_{\text{gb}}$.
Approximations were performed for every statisticaly significant signal of \kstar in different centrality and $p_T$ bins for each dataset. 
All approximations are presented on Fig. \ref{fig:minv}.

To ensure the validity of BG subtraction and approximation of \kstar signal FG/BG ratio were investigated. 
An example is shown on Fig. \ref{fig:minv_summary_example}. 
FG/BG ratios must approach 1 from above at high $M_{\text{inv}}$, and no statisticaly significant points below 1 should be present on FG/BG distribution. 
Signal of \kstar must be visible on FG/BG distribution (since \kstar signal is non-combinatorial) and background fit for \kstar signal, when divided by BG, must have a reasonable form in FG/BG.
